\documentclass[11pt,a4paper]{article}
\usepackage[T1]{fontenc}
\usepackage[utf8]{inputenc}
\usepackage{lmodern,microtype}
\usepackage[margin=26mm]{geometry}
\usepackage{amsmath,amssymb,amsthm,mathtools,booktabs}
\usepackage{xcolor}
\usepackage[colorlinks=true,allcolors=blue!50!black,breaklinks=true]{hyperref}
\usepackage{enumitem}
\setlist{nosep}
\setlength{\emergencystretch}{2em}
\newtheorem{theorem}{Theorem}[section]
\newtheorem{proposition}[theorem]{Proposition}
\newtheorem{corollary}[theorem]{Corollary}
\newtheorem{lemma}[theorem]{Lemma}
\theoremstyle{remark}
\newtheorem{remark}[theorem]{Remark}
\DeclareMathOperator{\per}{per}
\DeclareMathOperator{\Tr}{Tr}
\DeclareMathOperator{\diag}{diag}
\DeclareMathOperator{\rank}{rank}
\newcommand{\ket}[1]{|#1\rangle}
\newcommand{\bra}[1]{\langle#1|}
\newcommand{\ip}[2]{\langle #1,#2\rangle}
\newcommand{\T}{\mathbb T}
\newcommand{\E}{\mathcal E}
\newcommand{\Id}{\mathrm{Id}}
\title{Phase calibration unlocks adaptive tests\\for renewed quantum noise}
\author{Lluis Eriksson\\\small Independent researcher}
\date{15 October 2026\\\small Working manuscript}
\begin{document}
\maketitle
\begin{abstract}
Can interventions distinguish an unknown phase that is frozen throughout an
experiment from phase noise renewed at every use? We study ordinary forward
queries, arbitrary quantum memory, and exactly zero false positives for every
frozen diagonal unitary. For any Schur channel with correlation matrix \(C\),
the optimal parallel failure probability is
\(\min_{|m|=n}\per(C[m])/n!\), where \(C[m]\) repeats rows and columns according
to the occupation vector \(m\). The same formula is the exact optimum over
all causal strategies when the renewed channel has an unknown common phase
carrier and failure is evaluated in the worst case over that carrier.
Calibration changes this conclusion: for qubit dephasing we give a rational
five-query certificate that beats every parallel strategy for
\(0.5331424052\ldots<\eta<1\). At \(\eta=0.7\), failure decreases from
\(0.46603\) to \(0.4602957928558\). We prove physical realizability of the
certificate, obtain exact causal limits at complete dephasing, and show that
adaptation cannot improve the asymptotic failure exponent. The separation
concerns finite-query performance under a calibrated statistical promise;
it neither quantifies the hardware cost of a phase reference nor certifies
a quantum environment.
\end{abstract}

\section{Question, contribution, and scope}

A random phase can be sampled once and retained, or sampled anew at each
interaction. These mechanisms can have identical one-use channels and
different multi-time behavior. The operational role of interventions and
multi-time descriptions is well established \cite{pollock}. Our question is
more specific: how much can coherent control between uses improve
\emph{unambiguous exclusion of all frozen phases}?

There are two distinct promises for the renewed alternative. With a
\emph{calibrated carrier}, its channel is known. With an \emph{uncalibrated
carrier}, the channel is known only up to a fixed diagonal unitary applied
at every slot, and a test must perform well for every such carrier. This
distinction determines whether a finite-query adaptive advantage can occur.
Theorem~\ref{thm:minimax} rules out any advantage in the uncalibrated minimax
problem, in every dimension and at every query count.
Theorem~\ref{thm:separation} supplies a calibrated counterpoint using
commuting qubit noise. Here and below, \emph{adaptive} includes coherent
intermediate controls and memory; it does not require a recorded intermediate
measurement.

The permanent itself is not new to unambiguous comparison: Chefles,
Andersson, and Jex obtained \(\per(\mathrm{Gram})/N!\) as a symmetric-projector
failure probability \cite[Eq.~(4.10)]{chefles}. Our parallel theorem optimizes
the input, including arbitrary ancillas, for a channel problem indexed by
occupation types. Its extension to all causal tests under a common
unknown carrier is the principal structural result.
Likewise, causal Gram reconstruction is a specialization of established
query methods, notably Block Sets \cite[Definition~2, Theorems~2--3]{grillo},
semidefinite query characterizations \cite{barnum}, and quantum networks
\cite{combs}. We provide the specialized construction because a numerical
matrix inequality alone would not prove a physical adaptive separation.

Earlier unitary-comparison work \cite{andersson,sedlak,hashimoto} provides
direct context, but uses different alternatives or access models.
In particular, the updated multi-query comparison theorem of
Hashimoto, Soeda, and Murao \cite{hashimoto} establishes parallel optimality
for comparison of two Haar-unitary blocks, including general testers.
Exact adaptive sensing limits for bosonic dephasing \cite{huang} and
adaptive discrimination of specified channels \cite{sieniawski} also
predate this work. We claim neither a general invention of adaptive
discrimination nor the first exact dephasing limit.

The author's earlier work concerns sampling static disorder without
interventions within one realization \cite{sampling}, polynomial query
support \cite{support}, and an exact two-query adaptive separation for
tetrahedral echoes \cite{tetrahedral}. The present contribution is the
all-\(d,n,C\) calibration dichotomy together with its certified qubit
separation. Unlike the tetrahedral example, the calibrated causal optimum
here remains open.

\section{Model and notation}

Fix \(d\geq2\) and \(n\geq1\). A query acts on a \(d\)-dimensional register
in a known basis. Under the null, every slot applies the same unknown
\[
 D_\theta=\diag(e^{i\theta_0},\ldots,e^{i\theta_{d-1}}),
 \qquad \theta\in\T^d.
\]
It may be sampled from any distribution once per run. The experiment
outputs either rejection of the frozen-phase null or an inconclusive
outcome. Its rejection probability must be exactly zero for \emph{every}
\(\theta\). We call the inconclusive probability under the alternative
the failure probability \(\beta\). This is a one-sided test; it does not
unambiguously certify the null.

The calibrated alternative consists of independent uses of the Schur channel
\begin{equation}\label{eq:schur}
 \E_C(\ket a\bra b)=C_{ab}\ket a\bra b,\qquad C\succeq0,\quad C_{aa}=1.
\end{equation}
A dilation sends \(\ket a\) to \(\ket a\ket{e_a}\), with
\(\ip{e_b}{e_a}=C_{ab}\), using a fresh environment at each slot.
For a classical random-phase interpretation one additionally needs
\[
 C\in\operatorname{conv}\{zz^\dagger:|z_a|=1\}.
\]
This extra condition holds for the qubit channels below, but does not
follow from positive semidefiniteness and unit diagonal in general
dimension \(d\geq4\) \cite{dykema}. Our general theorems concern Schur
channels; their classical phase-renewal interpretation has this restriction.

Allowed tests use ordinary forward queries, finite ancillary systems,
ideal controls, and a final two-outcome measurement. Intermediate
instruments and randomness may be purified and stored in memory. We do
not allow inverse queries, controlled bypass of the queried device,
environment access, or indefinite causal order.
A parallel test prepares all \(n\) inputs jointly, queries them without
intervening controls, and measures jointly. Write
\(\beta_{\rm par}(n,C)\) and \(\beta_{\rm causal}(n,C)\) for the infima over
parallel and all causal tests respectively.

A word \(x=(x_1,\ldots,x_n)\in\{0,\ldots,d-1\}^n\) has occupation
\(m(x)=(m_0,\ldots,m_{d-1})\).
Let
\[
 \Omega_m=\{x:m(x)=m\},\qquad N_m=|\Omega_m|
       =\frac{n!}{\prod_a m_a!},\qquad |m|=n.
\]
For a representative word \(x\in\Omega_m\), let \(C[m]\) be the \(n\times n\)
matrix with entries \(C_{x_i x_j}\); its permanent is independent of the
representative. Set
\begin{equation}\label{eq:F}
 F_m(C)=\frac{\per(C[m])}{n!}.
\end{equation}

\section{The exact parallel optimum}

\begin{theorem}[Optimized permanent law]\label{thm:parallel}
For every correlation matrix \(C\), every \(d\), and \(n\geq1\),
\begin{equation}\label{eq:parallel}
 \beta_{\rm par}(n,C)=\min_{|m|=n}F_m(C).
\end{equation}
An optimum uses no ancilla: prepare
\(\ket{D_m}=N_m^{-1/2}\sum_{x\in\Omega_m}\ket x\) for a minimizing type,
and use the inconclusive effect \(\ket{D_m}\bra{D_m}\).
\end{theorem}
\begin{proof}
Purify an arbitrary parallel probe and write it as
\(\ket\psi=\sum_x\ket x\ket{a_x}\), with
\(p_x=\|\ket{a_x}\|^2\) and \(\sum_xp_x=1\).
Its null output is
\[
 \sum_m e^{im\cdot\theta}\ket{\psi_m},
 \qquad \ket{\psi_m}=\sum_{x\in\Omega_m}\ket x\ket{a_x}.
\]
Distinct characters are linearly independent, and the nonzero
\(\psi_m\) are orthogonal. Thus the null span is their direct sum.
If \(A\) is the inconclusive effect, zero rejection forces it to be
identity on that span; since \(0\preceq A\preceq I\), it dominates
the corresponding orthogonal projector.

Put \(K=C^{\otimes n}\) and \(p_m=\sum_{x\in\Omega_m}p_x\).
The resulting lower bound on failure is
\begin{equation}\label{eq:parallel-energy}
 \sum_{m:p_m>0}\frac{1}{p_m}
     \sum_{x,y\in\Omega_m}p_xp_y K_{xy}.
\end{equation}
This expression retains arbitrary ancillas: the orthogonal word registers
make the two inner products in each projector expectation equal to
\(p_x\) and \(p_y\).
Each principal block \(K_m=K|_{\Omega_m}\) is positive semidefinite.
Consequently \(q\mapsto q^\dagger K_m q\) is convex on the real probability
simplex. Slot permutations preserve \(K_m\) and act transitively on
\(\Omega_m\); averaging \(q_x=p_x/p_m\) over those permutations gives the
uniform distribution without increasing the energy.

The uniform energy is the constant row sum of \(K_m\) divided by \(N_m\).
In the permanent of \(C[m]\), each distinct word of type \(m\) occurs
\(\prod_a m_a!\) times. Hence this energy is
\[
 \frac{\per(C[m])}{N_m\prod_a m_a!}=F_m(C).
\]
Equation~\eqref{eq:parallel-energy} is bounded below by
\(\sum_m p_m F_m(C)\), and thus by the smallest \(F_m(C)\).
The stated uniform-type probe and projector attain equality.
\end{proof}

\begin{remark}
The argument includes complex and singular \(C\). It does not assume
entrywise nonnegativity. Its permanent is the familiar Gram-permanent
quantity from state comparison \cite{chefles}; the additional statement
is global optimization over channel probes and ancillas.
\end{remark}

\section{Unknown carrier: a minimax theorem for all causal tests}

Suppose the alternative is known only up to a fixed common carrier
\(\phi\in\T^d\):
\[
 C^{(\phi)}=D_\phi C D_\phi^\dagger,\qquad
 \E_{C^{(\phi)}}=\mathcal U_\phi\circ\E_C,\quad
 \mathcal U_\phi(X)=D_\phi XD_\phi^\dagger.
\]
Conditional on \(\phi\), the \(n\) uses are independent. The carrier is
fixed throughout the run, rather than drawn anew at every slot.
Define
\[
 \beta_{\rm causal}^{\rm wc}(n,C)
 =\inf_{T:\,\text{zero null error}}\ \sup_{\phi\in\T^d}
                   \beta_T(C^{(\phi)}).
\]
The infimum includes strategies that attempt to learn \(\phi\) within
their \(n\) queries.

\begin{theorem}[Calibration-free minimax equality]\label{thm:minimax}
For every \(d,n,C\) as above,
\begin{equation}\label{eq:minimax}
 \beta_{\rm causal}^{\rm wc}(n,C)
 =\beta_{\rm par}(n,C)=\min_{|m|=n}F_m(C).
\end{equation}
Thus a uniform-type parallel test is minimax optimal among all causal tests.
\end{theorem}
\begin{proof}
A purified causal circuit has a phase-history expansion
\begin{equation}\label{eq:history}
 \ket{\Psi(z_1,\ldots,z_n)}
   =\sum_x\left(\prod_{j=1}^n z_{j,x_j}\right)\ket{v_x},
 \qquad |z_{j,a}|=1.
\end{equation}
Let \(u_m=\sum_{x\in\Omega_m}v_x\). Under the frozen null the output is
\(\sum_m e^{im\cdot\theta}u_m\), of norm one for every \(\theta\).
Haar integration and Parseval give
\begin{equation}\label{eq:parseval}
 \sum_m\|u_m\|^2=1.
\end{equation}
The vectors \(u_m\) need not be orthogonal.

Let \(A\) be the inconclusive effect. Zero rejection gives
\(A\Psi_\theta=\Psi_\theta\). Fourier extraction then implies
\(\sqrt A\,u_m=u_m\). Write \(a_x=\sqrt A\,v_x\), so that
\(\sum_{x\in\Omega_m}a_x=u_m\).
The alternative output, by the fresh-environment dilation of
\eqref{eq:schur}, has acceptance probability
\[
 \sum_{x,y}K_{xy}
 e^{i(m(x)-m(y))\cdot\phi}\ip{a_y}{a_x}.
\]
Averaging \(\phi\) deletes only cross-type terms. On one type block,
slot-permutation invariance implies
\(K_m\mathbf1=N_mF_m(C)\mathbf1\).
Decompose
\[
 a_x=\frac{u_m}{N_m}+b_x,\qquad
 \sum_{x\in\Omega_m}b_x=0.
\]
The mixed terms vanish by the constant row sum, and positivity of \(K_m\)
gives
\begin{equation}\label{eq:block}
 \sum_{x,y\in\Omega_m}K_{xy}\ip{a_y}{a_x}
 =F_m(C)\|u_m\|^2+
   \sum_{x,y\in\Omega_m}K_{xy}\ip{b_y}{b_x}
 \ \geq\ F_m(C)\|u_m\|^2.
\end{equation}
For complex \(C\), the last energy is nonnegative because it is the
contraction of two positive Gram matrices (equivalently, a sum of squared
tensor-product norms).
Sum \eqref{eq:block} and apply \eqref{eq:parseval}. The mean failure over
\(\phi\), hence its supremum, is at least \(\min_m F_m(C)\).
No interchange of infimum and supremum is required.
The uniform-type test in Theorem~\ref{thm:parallel} has the same failure
for every carrier, since \(m(x)=m(y)\) throughout its support.
It attains the lower bound.
\end{proof}

\begin{corollary}[Complete dephasing]\label{cor:complete}
For \(C=I_d\), even with a calibrated carrier,
\begin{equation}\label{eq:complete}
 \beta_{\rm causal}(n,I_d)
 =\beta_{\rm par}(n,I_d)
 =\frac{1}{\max_{|m|=n}\binom{n}{m_0,\ldots,m_{d-1}}}.
\end{equation}
The maximum occurs when occupations differ by at most one.
\end{corollary}
\begin{proof}
The channel is invariant under carrier shifts, so its minimax and
calibrated problems coincide. Here \(K_m\) is identity and \(F_m=1/N_m\).
If \(m_a\geq m_b+2\), moving one count from \(a\) to \(b\) decreases
\(\prod_a m_a!\); iterating proves the balancing assertion.
\end{proof}

\section{Qubit formulas and the asymptotic limit}

For qubits take a known zero carrier and
\begin{equation}\label{eq:eta}
 C_\eta=\begin{pmatrix}1&\eta\\ \eta&1\end{pmatrix},
 \quad 0\leq\eta\leq1,\qquad
 \E_\eta=p\Id+q\mathcal Z,\quad
 p=\frac{1+\eta}{2},\quad q=\frac{1-\eta}{2},
\end{equation}
where \(\mathcal Z(X)=ZXZ\) and \(Z=\diag(1,-1)\).
For a known nonzero carrier, a compensating control gives this model
while preserving the composite frozen null.

\begin{proposition}[Balanced qubit formula]\label{prop:qubit}
With \(w=\lfloor n/2\rfloor\),
\begin{equation}\label{eq:qubit}
 \beta_{\rm par}(n,\eta)
 =\frac{\displaystyle\sum_{j=0}^{w}
            \binom{w}{j}\binom{n-w}{j}\eta^{2j}}{\binom{n}{w}}.
\end{equation}
\end{proposition}
\begin{proof}
In a weight-\(w\) block, a fixed binary word has
\(\binom wj\binom{n-w}j\) words at Hamming distance \(2j\), giving the
displayed value for that weight.
Let \(J_w\) have probabilities
\(\pi_w(j)=\binom wj\binom{n-w}j/\binom nw\).
For \(w<(n-1)/2\), the likelihood ratio on the common support is
\[
 \frac{\pi_{w+1}(j)}{\pi_w(j)}
 =\frac{(w+1)^2(n-w-j)}{(n-w)^2(w+1-j)},
\]
which increases with \(j\); any new support lies at the upper end.
Thus \(J_{w+1}\) stochastically dominates \(J_w\). Since
\(\eta^{2j}\) decreases in \(j\), the block value decreases toward the
balanced weight. Symmetry handles the other half. At \(\eta=0\) the
value is \(1/\binom nw\), and at \(\eta=1\) all values are one.
\end{proof}

\begin{proposition}[No improvement of the failure exponent]\label{prop:exp}
For \(n\geq1\) and \(0\leq\eta\leq1\),
\begin{equation}\label{eq:expbounds}
 p^n+q^n\ \leq\ \beta_{\rm causal}(n,\eta)
 \ \leq\ \beta_{\rm par}(n,\eta)\ \leq\ (n+1)p^n.
\end{equation}
Consequently, for fixed \(0\leq\eta<1\), both optimal failure exponents
are \(-\log((1+\eta)/2)\).
\end{proposition}
\begin{proof}
The Kraus trajectories consisting entirely of \(I\), or entirely of
\(Z\), have probabilities \(p^n\) and \(q^n\). Each is a frozen null and
must be inconclusive with certainty, for any causal protocol.
This proves the lower bound. In \eqref{eq:qubit}, the numerator is at most
\((1+\eta)^n\): it is a sub-sum of the product
\((1+\eta)^w(1+\eta)^{n-w}\) with nonnegative terms.
Also \(\binom n{\lfloor n/2\rfloor}\geq2^n/(n+1)\).
Taking logarithms and dividing by \(n\) proves the limit.
\end{proof}

\section{Causal certificates and physical realization}

For completeness we specialize known query-decomposition methods
\cite{grillo,barnum,combs} to the present continuous-phase setting.
Index level \(k\) by words of length \(k\), with the last symbol appended
on the right. Let \(V\) be the \(d^n\times\binom{n+d-1}{d-1}\) matrix whose
columns are the type indicators \(\mathbf1_{\Omega_m}\).

\begin{lemma}[Causal Gram realization]\label{lem:gram}
The following conditions characterize causal rejection statistics:
\begin{align}
 &Q_0=[1],\qquad Q_k\succeq0,\label{eq:gram1}\\
 &(Q_k)_{xa,yb}=0\quad(a\ne b),\qquad
   \sum_a(Q_k)_{xa,ya}=(Q_{k-1})_{x,y},\label{eq:gram2}\\
 &0\preceq R\preceq Q_n,\qquad RV=0.\label{eq:gram3}
\end{align}
The rejection probability under \(\E_C^{\otimes n}\) is
\begin{equation}\label{eq:objective}
 s(C)=\Tr(C^{\otimes n}R).
\end{equation}
Every feasible collection has a realization using a \(d\)-dimensional
query register and memory of dimension at most \(d^{n-1}\).
\end{lemma}
\begin{proof}
Expand the purified output immediately after each query into its
phase histories, with coefficient vectors \(v_x^{(k)}\), and set
\((Q_k)_{x,y}=\ip{v_x^{(k)}}{v_y^{(k)}}\).
Projection onto the last queried symbol makes distinct last-symbol
branches orthogonal. Summing these projections preserves the previous
inner products, proving \eqref{eq:gram1}--\eqref{eq:gram2}.
For a final rejection effect \(M\), set
\(R_{x,y}=\ip{v_x^{(n)}}{M v_y^{(n)}}\).
Then \(0\preceq R\preceq Q_n\). Fourier extraction of the null characters
shows that zero rejection is equivalent to annihilating all type sums;
by positivity this is \(RV=0\).
The channel dilation contracts two histories with
\((C^{\otimes n})_{xy}\), giving \eqref{eq:objective}, including for
complex \(C\).

Conversely, factor the same-last-symbol blocks of \(Q_k\) as
\(B_a^\dagger B_a\). If previous coefficient vectors have Gram
\(Q_{k-1}\), define on their span
\[
 v_x^{(k-1)}\longmapsto \sum_a\ket a\otimes B_a\ket x.
\]
The Gram recursion makes this map well defined and isometric, including
when the previous vectors are linearly dependent. Extend the isometry
to a unitary after completing the ambient spaces. Applying the next
diagonal query produces exactly the required next coefficient vectors.
This constructs all levels.

If \(S\) has the terminal vectors as columns, \(S^\dagger S=Q_n\).
The inequality \(0\preceq R\preceq S^\dagger S\) ensures that \(R\)
annihilates \(\ker S\). The effect
\(M=(S^+)^\dagger R S^+\) on \(\operatorname{ran}S\), extended by zero,
satisfies \(0\preceq M\preceq I\) and \(S^\dagger M S=R\).
Each same-symbol block has size at most \(d^{k-1}\).
Embedding all levels in a space of dimension \(d^n\) proves the stated
sufficient memory bound.
\end{proof}

This lemma is an exact physical characterization, not a relaxation
based only on polynomial normalization. It supplies a convenient
specialization of the earlier reconstruction principles; we do not
claim their general equivalence as new. The sufficient memory bound
also says nothing about efficient compilation into a restricted gate set.

\section{A certified calibrated separation}

At \(n=5\), \eqref{eq:qubit} becomes
\begin{equation}\label{eq:fivepar}
 \beta_{\rm par}(5,\eta)=\frac{1+6\eta^2+3\eta^4}{10}.
\end{equation}
The supplementary files \texttt{phase-certificate-Q.json} and
\texttt{phase-certificate-R.json} specify rational \(32\times32\)
matrices, in binary lexicographic order. They define a single fixed
strategy, independent of \(\eta\).

\begin{theorem}[Five-query calibrated advantage]\label{thm:separation}
There is a causal five-query test with exactly zero rejection of every
frozen phase, whose rejection probability is
\begin{align}
 s_*(\eta)={}&\frac{1086981}{1250000}
 +\frac{3225799}{50000000}\eta
 -\frac{3784393}{6250000}\eta^2 \notag\\
 &-\frac{585391}{25000000}\eta^3
 -\frac{103157}{390625}\eta^4
 -\frac{2055017}{50000000}\eta^5.\label{eq:success}
\end{align}
It strictly improves on every parallel test for
\(\eta_*<\eta<1\), where \(\eta_*\) is the unique real root of
\[
 g(\eta)=2055017\eta^3-1795904\eta^2+3225799\eta-1520760
\]
and \(0.53314<\eta_*<0.53315\).
At \(\eta=7/10\), its failure and improvement are respectively
\begin{equation}\label{eq:exactgap}
 1-s_*(7/10)=\frac{2301478964279}{5000000000000},\qquad
 \beta_{\rm par}(5,7/10)-(1-s_*(7/10))
 =\frac{28671035721}{5000000000000}.
\end{equation}
\end{theorem}
\begin{proof}[Proof with a finite rational certificate]
Reduce \(Q=Q_5\) recursively by
\((Q_{k-1})_{xy}=(Q_k)_{x0,y0}+(Q_k)_{x1,y1}\).
The supplied entries satisfy all off-last-symbol zeros and \(Q_0=[1]\).
Exact rational Schur-complement elimination establishes
\[
 Q\succ0,\quad Q-R\succ0,\quad R\succeq0,\qquad
 \rank Q=\rank(Q-R)=32,\quad\rank R=26.
\]
The six type-indicator columns are annihilated exactly by \(R\).
Lemma~\ref{lem:gram} therefore supplies a physical test with one query
qubit and at most a 16-dimensional memory.

Summing \(R_{xy}\) over each Hamming distance gives the six coefficients
in \eqref{eq:success}. Subtracting its failure from \eqref{eq:fivepar}
gives the exact identity
\begin{equation}\label{eq:factor}
 \beta_{\rm par}(5,\eta)-(1-s_*(\eta))
       =\frac{(1-\eta^2)g(\eta)}{50000000}.
\end{equation}
The derivative of \(g\) has positive leading coefficient and negative
discriminant \(-66647776694132\), so \(g\) increases on the real line.
Exact substitution at \(0.53314\) and \(0.53315\) brackets its root.
Equations~\eqref{eq:factor} and \eqref{eq:exactgap} follow.

The included verifier replays all checks in exact rational arithmetic.
No semidefinite optimizer is needed.
Positive semidefiniteness is checked by successive Schur complements:
a negative pivot fails, and a zero pivot is allowed only if its remaining
row is identically zero. This yields an exact finite-arithmetic proof
for the supplied entries.
\end{proof}

\begin{center}
\begin{tabular}{lrr}
\toprule
Test at \(\eta=0.7\) & Failure & Detection\\
\midrule
Optimal parallel & \(0.4660300000000\) & \(0.5339700000000\)\\
Certified causal & \(0.4602957928558\) & \(0.5397042071442\)\\
\bottomrule
\end{tabular}
\end{center}

The improvement is about \(0.5734\) percentage points of detection.
It is mathematically strict but modest. The theorem does not identify
the optimal causal test, prove that five queries are necessary, or
improve the exponent in Proposition~\ref{prop:exp}. In particular, a
floating-point optimizer value is not an additional exact theorem.
At \(\eta=0\), the exact optimum is \(1/\binom52=0.1\) by
Corollary~\ref{cor:complete}; the displayed certificate is not optimal
there. At \(\eta=1\), the alternative is itself a frozen unitary.

\section{Equal marginals and perturbations}

\begin{proposition}[Matched one-use channels]\label{prop:matching}
For \(0\leq\eta<1\), requiring zero rejection against every frozen qubit
phase law with first moment \(\eta\) is equivalent to requiring zero
rejection for each deterministic frozen phase.
\end{proposition}
\begin{proof}
Fix a point \(z\) on the unit circle. Choose
\(0<\epsilon\leq(1-\eta)/2\).
The residual mean \((\eta-\epsilon z)/(1-\epsilon)\) lies in the unit
disk, because \(|\eta-\epsilon z|\leq\eta+\epsilon\leq1-\epsilon\).
It is a convex combination of two circle points. Thus a phase law
of mean \(\eta\) can assign positive mass \(\epsilon\) to any prescribed
\(z\). Zero rejection averaged over every such law, together with
nonnegativity, forces zero rejection at that point. The converse is
immediate.
\end{proof}

In particular, sampling \(I\) or \(Z\) once with probabilities \(p,q\)
and retaining it has the same one-use channel as independently
sampling them at every slot. The proposition concerns \emph{all}
compatible frozen laws, rather than a known two-point null law alone.
Rejection excludes the specified frozen-phase null within the promise;
it does not establish universal Markovianity, an intrinsically quantum
bath, or a unique physical noise mechanism.

\begin{proposition}[Qubit null perturbations]\label{prop:robust}
For any zero-null-error qubit test, suppose the relative phase at slot
\(j\) is \(\theta+\delta_j\), with \(|\delta_j|\leq\pi\).
Its false-positive probability is at most
\begin{equation}\label{eq:robust}
 \min\left\{1,\left(\frac12\sum_{j=1}^n|\delta_j|\right)^2\right\}.
\end{equation}
A common phase change across all slots remains an exact null.
\end{proposition}
\begin{proof}
After removing a global scalar, the operator-norm change of one qubit
phase query is \(2|\sin(\delta_j/4)|\leq|\delta_j|/2\).
Telescoping the purified isometric circuit bounds the final vector
change by the sum. The square root of the rejection effect annihilates
the ideal null vector and is a contraction. Squaring gives
\eqref{eq:robust}.
\end{proof}

This is a bound on a fixed test's false positives. It does not prove
superiority to the best parallel test allowed a positive false-positive
level. Translating time errors into \(\delta_j\) requires a frequency
or phase-width bound. The result is not a claim of clock robustness
without a physical scale, and it does not include arbitrary state
preparation, measurement, or memory errors.

\section{Reproducibility and remaining questions}

The certificate matrices were obtained by numerical semidefinite search,
interior mixing, rational rounding, and exact constraint repair.
The final proof depends on the rational entries and their verifier,
not on optimizer tolerances. The two matrix files have SHA-256 hashes
\begin{quote}\small\ttfamily
Q: 18ecfa37bc84e3b0ca40efdc9dbdda6b87442479d76a17f4f7ff04563f088463\\
R: 8a3bc80cdc93463652359b7145a095ac528c1469486173fb3f3cc3ad8e4000be
\end{quote}
The standard-library verifier checks causal constraints, positivity,
null annihilation, coefficient aggregation, and the exact separation.
A separate numerical reconstruction checks the circuit against frozen
phases, independent phase histories, and direct Kraus evolution; its
largest probability discrepancy was below \(10^{-14}\).
These numerical checks are diagnostics, not proofs of continuum claims.

A separate GPT-6 Astra review context checked the candidate proofs,
reimplemented rational verification, reconstructed the circuit, and
tested the minimax argument. That review shares the model provider and
was coordinated with the authoring process; it is neither blind review
nor human peer review. Its report and scope are included with the
supplement. AI systems assisted literature triage, mathematical
development, coding, drafting, and review. No formal proof-assistant
verification is claimed.

The literature search renewed indexes from frontier-AI mathematics
collections and the author's repositories, followed by targeted reading
of the closest quantum-information sources. Index coverage is not
full-text mathematical verification. The accompanying search audit
records that distinction; the absence of a located identical result is
not a guarantee of priority.

Three questions remain especially concrete. First, what is the exact
calibrated causal optimum, and what is the smallest query count admitting
a strict advantage? Second, can the five-query certificate be replaced
by a low-memory analytic circuit that explains the beneficial
interference between occupation sectors? Third, does a useful advantage
persist after jointly optimizing both test classes at a positive
false-positive allowance and with realistic control costs?
The present results settle the uncalibrated minimax problem and exhibit
a finite calibrated separation, while leaving these stronger claims open.

\begin{thebibliography}{99}
\bibitem{chefles}
A. Chefles, E. Andersson, and I. Jex,
\emph{Unambiguous comparison of the states of multiple quantum systems},
J. Phys. A \textbf{37}, 7315 (2004).
\href{https://arxiv.org/abs/quant-ph/0402125}{arXiv:quant-ph/0402125}.
\bibitem{grillo}
S. A. Grillo and F. L. Marquezino,
\emph{Quantum Query as a State Decomposition},
Theoretical Computer Science (2018),
\href{https://doi.org/10.1016/j.tcs.2018.03.017}{doi:10.1016/j.tcs.2018.03.017}.
\href{https://arxiv.org/abs/1602.07716v2}{arXiv:1602.07716v2}.
\bibitem{barnum}
H. Barnum,
\emph{Semidefinite programming characterization and spectral adversary
method for quantum complexity with noncommuting unitary queries} (2007).
\href{https://arxiv.org/abs/quant-ph/0703141}{arXiv:quant-ph/0703141}.
\bibitem{combs}
G. Chiribella, G. M. D'Ariano, and P. Perinotti,
\emph{Theoretical framework for quantum networks},
Phys. Rev. A \textbf{80}, 022339 (2009).
\href{https://arxiv.org/abs/0904.4483}{arXiv:0904.4483}.
\bibitem{andersson}
E. Andersson, I. Jex, and S. M. Barnett,
\emph{Comparison of unitary transforms},
J. Phys. A \textbf{36} (2003),
\href{https://doi.org/10.1088/0305-4470/36/9/310}{doi:10.1088/0305-4470/36/9/310}.
\href{https://arxiv.org/abs/quant-ph/0208153}{arXiv:quant-ph/0208153}.
\bibitem{sedlak}
M. Sedl\'ak and M. Ziman,
\emph{Unambiguous comparison of unitary channels},
Phys. Rev. A \textbf{79}, 012303 (2009).
\href{https://arxiv.org/abs/0809.4401}{arXiv:0809.4401}.
\bibitem{hashimoto}
Y. Hashimoto, A. Soeda, and M. Murao,
\emph{Comparison of unknown unitary channels with multiple queries},
\href{https://arxiv.org/abs/2208.12519v2}{arXiv:2208.12519v2}
(version of 7 September 2026).
\bibitem{huang}
Z. Huang, L. Lami, and M. M. Wilde,
\emph{Exact Quantum Sensing Limits for Bosonic Dephasing Channels},
PRX Quantum \textbf{5}, 020354 (2024).
\href{https://arxiv.org/abs/2402.05793}{arXiv:2402.05793}.
\bibitem{sieniawski}
S. Sieniawski and R. Demkowicz-Dobrza\'nski,
\emph{Adaptive quantum channel discrimination using methods of quantum
metrology}, New J. Phys. \textbf{28}, 024502 (2026).
\href{https://arxiv.org/abs/2510.15506}{arXiv:2510.15506}.
\bibitem{pollock}
F. A. Pollock, C. Rodr\'iguez-Rosario, T. Frauenheim,
M. Paternostro, and K. Modi,
\emph{Operational Markov condition for quantum processes},
Phys. Rev. Lett. \textbf{120}, 040405 (2018).
\href{https://arxiv.org/abs/1801.09811}{arXiv:1801.09811}.
\bibitem{dykema}
K. Dykema and K. Juschenko,
\emph{Matrices of unitary moments} (2009).
\href{https://arxiv.org/abs/0901.0288}{arXiv:0901.0288}.
\bibitem{sampling}
L. Eriksson, \emph{The arithmetic of sampling static quantum disorder:
exact aliases, finite witnesses, and clock robustness}, working paper.
\href{https://airr.science/papers/ARR-2026-7PEKNEDVH38HV8PX/versions/v1/}
{AIRR ARR-2026-7PEKNEDVH38HV8PX, v1 (2026)}.
\bibitem{support}
L. Eriksson, \emph{Algebraic Query Support for Unitary Oracles: Exact
Hilbert Laws, Harmonic Spectra, and General-Tester Bounds}, working paper.
\href{https://airr.science/papers/ARR-2026-6WX2JF38WE87GB2M/versions/v1/}
{AIRR ARR-2026-6WX2JF38WE87GB2M, v1 (2026)}.
\bibitem{tetrahedral}
L. Eriksson, \emph{Two-Query Chirality in Tetrahedral Quantum Echoes:
Exact Parallel Readout, a Nine-Dimensional Query Defect, and a Certified
Adaptive Advantage}, working paper.
\href{https://airr.science/papers/ARR-2026-5KS70GV7KK9DYA69/versions/v1/}
{AIRR ARR-2026-5KS70GV7KK9DYA69, v1 (2026)}.
\end{thebibliography}
\end{document}
